% paper75-copilot-deduction-assist.tex
%   CopilotDeductionAssist — Trust-Calibrated Deduction Hints
%   Paper 75 of the Judgment Geometry series.
% Compile: pdflatex paper75-copilot-deduction-assist
\documentclass[11pt]{article}
\usepackage{lmodern}
% jugeo-common.tex — shared preamble for all Judgment Geometry papers
% Include via: % jugeo-common.tex — shared preamble for all Judgment Geometry papers
% Include via: % jugeo-common.tex — shared preamble for all Judgment Geometry papers
% Include via: \input{jugeo-common}

% ── Packages ────────────────────────────────────────────────────────────
\usepackage{amsmath,amssymb,amsthm,mathtools}
\usepackage{tikz-cd}
\usepackage{listings}
\usepackage{xcolor}
\usepackage{xspace}
\usepackage{booktabs}
\usepackage{multirow}
\usepackage{enumitem}
\usepackage[expansion=false]{microtype}
\usepackage{hyperref}
\usepackage{cleveref}
% \usepackage{thmtools}  % disabled: not available in this TeX installation
\usepackage{float}
\usepackage{caption}
\usepackage{subcaption}
\usepackage{tabularx}
\usepackage{stmaryrd}       % \llbracket, \rrbracket
\usepackage{bbm}            % \mathbbm{1}

% ── Page geometry ───────────────────────────────────────────────────────
\usepackage[margin=1in]{geometry}

% ── Theorem environments ────────────────────────────────────────────────
\theoremstyle{plain}
\newtheorem{theorem}{Theorem}[section]
\newtheorem{lemma}[theorem]{Lemma}
\newtheorem{proposition}[theorem]{Proposition}
\newtheorem{corollary}[theorem]{Corollary}

\theoremstyle{definition}
\newtheorem{definition}[theorem]{Definition}
\newtheorem{example}[theorem]{Example}
\newtheorem{construction}[theorem]{Construction}

\theoremstyle{remark}
\newtheorem{remark}[theorem]{Remark}
\newtheorem{notation}[theorem]{Notation}

% ── Python listings style ──────────────────────────────────────────────
\definecolor{jugeoBlue}{HTML}{2563EB}
\definecolor{jugeoGreen}{HTML}{059669}
\definecolor{jugeoRed}{HTML}{DC2626}
\definecolor{jugeoPurple}{HTML}{7C3AED}
\definecolor{jugeoGray}{HTML}{6B7280}
\definecolor{jugeoBg}{HTML}{F8FAFC}

\lstdefinestyle{jugeo-python}{
  language=Python,
  basicstyle=\ttfamily\small,
  keywordstyle=\color{jugeoBlue}\bfseries,
  stringstyle=\color{jugeoGreen},
  commentstyle=\color{jugeoGray}\itshape,
  numberstyle=\tiny\color{jugeoGray},
  backgroundcolor=\color{jugeoBg},
  numbers=left,
  numbersep=6pt,
  frame=single,
  framerule=0.4pt,
  rulecolor=\color{jugeoGray!40},
  breaklines=true,
  breakatwhitespace=true,
  showstringspaces=false,
  tabsize=4,
  xleftmargin=1.5em,
  framexleftmargin=1.5em,
  morekeywords={dataclass,frozen,field,Enum,IntEnum,auto,Any,
                Mapping,Sequence,Optional,match,case},
  literate={→}{{$\to$}}1 {←}{{$\leftarrow$}}1
           {≤}{{$\leq$}}1 {≥}{{$\geq$}}1 {≠}{{$\neq$}}1
           {∈}{{$\in$}}1 {∀}{{$\forall$}}1 {∃}{{$\exists$}}1
           {⊕}{{$\oplus$}}1 {⊖}{{$\ominus$}}1,
  escapeinside={(*@}{@*)},
}
\lstset{style=jugeo-python}

\lstdefinestyle{jugeo-output}{
  basicstyle=\ttfamily\small,
  backgroundcolor=\color{jugeoBg},
  frame=single,
  framerule=0.4pt,
  rulecolor=\color{jugeoGray!40},
  breaklines=true,
  showstringspaces=false,
  xleftmargin=1.5em,
  framexleftmargin=0.5em,
  numbers=none,
}

% ── JuGeo-specific macros ──────────────────────────────────────────────

% Judgment 8-tuple
\newcommand{\J}{\mathcal{J}}
\newcommand{\Jud}[8]{(#1,\,#2,\,#3,\,#4,\,#5,\,#6,\,#7,\,#8)}
\newcommand{\JudShort}[1]{\J_{#1}}

% Components
\newcommand{\coord}{\mathsf{c}}            % coordinate
\newcommand{\prop}{\varphi}                 % proposition
\newcommand{\carrier}{\mathcal{A}}          % carrier type
\newcommand{\evid}{\mathcal{E}}             % evidence bundle
\newcommand{\oblig}{\mathcal{O}}            % obligations
\newcommand{\obstr}{\mathcal{B}}            % obstructions
\newcommand{\trust}{\mathcal{T}}            % trust annotation
\newcommand{\prov}{\Pi}                     % provenance

% Trust algebra
\newcommand{\Talg}{\mathfrak{T}}            % trust ordered algebra
\newcommand{\Eadm}{\mathcal{E}_{\mathrm{adm}}}
\newcommand{\tleq}{\preceq}                 % trust ordering
\newcommand{\tjoin}{\oplus}                 % conservative join
\newcommand{\tatten}{\ominus}               % attenuation
\newcommand{\tpromote}{\uparrow_\pi}        % promotion
\newcommand{\tdemote}{\downarrow_\chi}       % demotion

% Trust tiers
\newcommand{\Tcontra}{\ensuremath{\mathsf{CONTRADICTED}}}
\newcommand{\Tunver}{\ensuremath{\mathsf{UNVERIFIED}}}
\newcommand{\Tcopilot}{\ensuremath{\mathsf{COPILOT}}}
\newcommand{\Toracle}{\ensuremath{\mathsf{ORACLE}}}
\newcommand{\Truntime}{\ensuremath{\mathsf{RUNTIME}}}
\newcommand{\Tsolver}{\ensuremath{\mathsf{SOLVER}}}
\newcommand{\Tproof}{\ensuremath{\mathsf{PROOF}}}

% Site and geometry
\newcommand{\Site}{\mathcal{S}}             % semantic site
\newcommand{\Cov}{\mathrm{Cov}}            % covering family
\newcommand{\Ob}{\mathrm{Ob}}              % objects
\newcommand{\Mor}{\mathrm{Mor}}            % morphisms
\newcommand{\res}{\mathrm{res}}            % restriction
\newcommand{\Cech}{\check{C}}              % Čech complex
\newcommand{\Hcoh}[1]{H^{#1}}             % cohomology group

% Descent
\newcommand{\descent}{\mathrm{desc}}
\newcommand{\glue}{\mathrm{glue}}
\newcommand{\overlap}[2]{#1 \cap #2}

% Semantic moves
\newcommand{\VERIFY}{\textsc{Verify}}
\newcommand{\CONSTRUCT}{\textsc{Construct}}
\newcommand{\NORMALIZE}{\textsc{Normalize}}
\newcommand{\GLUE}{\textsc{Glue}}
\newcommand{\RESTRICT}{\textsc{Restrict}}
\newcommand{\TRANSPORT}{\textsc{Transport}}
\newcommand{\REPAIR}{\textsc{Repair}}
\newcommand{\PROMOTE}{\textsc{Promote}}
\newcommand{\CHALLENGE}{\textsc{Challenge}}
\newcommand{\ESCALATE}{\textsc{Escalate}}

% Fragments
\newcommand{\QFLIA}{\texttt{QF\_LIA}}
\newcommand{\QFLRA}{\texttt{QF\_LRA}}
\newcommand{\QFBV}{\texttt{QF\_BV}}
\newcommand{\QFUF}{\texttt{QF\_UF}}

% Misc
\newcommand{\Python}{\textsc{Python}}
\newcommand{\JuGeo}{\textsc{JuGeo}}
\newcommand{\LEAN}{\textsc{Lean}\,4}
\newcommand{\Fstar}{F$^{\star}$}
\newcommand{\Coq}{\textsc{Coq}}
\newcommand{\Dafny}{\textsc{Dafny}}

% Common shortcuts
\newcommand{\ie}{i.e.\@\xspace}
\newcommand{\eg}{e.g.\@\xspace}
\newcommand{\cf}{cf.\@\xspace}
\newcommand{\wrt}{w.r.t.\@\xspace}

% Evaluation shortcuts
\newcommand{\numcases}{300}
\newcommand{\numspec}{100}
\newcommand{\numeq}{100}
\newcommand{\numbug}{100}
\newcommand{\medtime}{6.5\,ms}
\newcommand{\totaltime}{1.949\,s}


% ── Packages ────────────────────────────────────────────────────────────
\usepackage{amsmath,amssymb,amsthm,mathtools}
\usepackage{tikz-cd}
\usepackage{listings}
\usepackage{xcolor}
\usepackage{xspace}
\usepackage{booktabs}
\usepackage{multirow}
\usepackage{enumitem}
\usepackage[expansion=false]{microtype}
\usepackage{hyperref}
\usepackage{cleveref}
% \usepackage{thmtools}  % disabled: not available in this TeX installation
\usepackage{float}
\usepackage{caption}
\usepackage{subcaption}
\usepackage{tabularx}
\usepackage{stmaryrd}       % \llbracket, \rrbracket
\usepackage{bbm}            % \mathbbm{1}

% ── Page geometry ───────────────────────────────────────────────────────
\usepackage[margin=1in]{geometry}

% ── Theorem environments ────────────────────────────────────────────────
\theoremstyle{plain}
\newtheorem{theorem}{Theorem}[section]
\newtheorem{lemma}[theorem]{Lemma}
\newtheorem{proposition}[theorem]{Proposition}
\newtheorem{corollary}[theorem]{Corollary}

\theoremstyle{definition}
\newtheorem{definition}[theorem]{Definition}
\newtheorem{example}[theorem]{Example}
\newtheorem{construction}[theorem]{Construction}

\theoremstyle{remark}
\newtheorem{remark}[theorem]{Remark}
\newtheorem{notation}[theorem]{Notation}

% ── Python listings style ──────────────────────────────────────────────
\definecolor{jugeoBlue}{HTML}{2563EB}
\definecolor{jugeoGreen}{HTML}{059669}
\definecolor{jugeoRed}{HTML}{DC2626}
\definecolor{jugeoPurple}{HTML}{7C3AED}
\definecolor{jugeoGray}{HTML}{6B7280}
\definecolor{jugeoBg}{HTML}{F8FAFC}

\lstdefinestyle{jugeo-python}{
  language=Python,
  basicstyle=\ttfamily\small,
  keywordstyle=\color{jugeoBlue}\bfseries,
  stringstyle=\color{jugeoGreen},
  commentstyle=\color{jugeoGray}\itshape,
  numberstyle=\tiny\color{jugeoGray},
  backgroundcolor=\color{jugeoBg},
  numbers=left,
  numbersep=6pt,
  frame=single,
  framerule=0.4pt,
  rulecolor=\color{jugeoGray!40},
  breaklines=true,
  breakatwhitespace=true,
  showstringspaces=false,
  tabsize=4,
  xleftmargin=1.5em,
  framexleftmargin=1.5em,
  morekeywords={dataclass,frozen,field,Enum,IntEnum,auto,Any,
                Mapping,Sequence,Optional,match,case},
  literate={→}{{$\to$}}1 {←}{{$\leftarrow$}}1
           {≤}{{$\leq$}}1 {≥}{{$\geq$}}1 {≠}{{$\neq$}}1
           {∈}{{$\in$}}1 {∀}{{$\forall$}}1 {∃}{{$\exists$}}1
           {⊕}{{$\oplus$}}1 {⊖}{{$\ominus$}}1,
  escapeinside={(*@}{@*)},
}
\lstset{style=jugeo-python}

\lstdefinestyle{jugeo-output}{
  basicstyle=\ttfamily\small,
  backgroundcolor=\color{jugeoBg},
  frame=single,
  framerule=0.4pt,
  rulecolor=\color{jugeoGray!40},
  breaklines=true,
  showstringspaces=false,
  xleftmargin=1.5em,
  framexleftmargin=0.5em,
  numbers=none,
}

% ── JuGeo-specific macros ──────────────────────────────────────────────

% Judgment 8-tuple
\newcommand{\J}{\mathcal{J}}
\newcommand{\Jud}[8]{(#1,\,#2,\,#3,\,#4,\,#5,\,#6,\,#7,\,#8)}
\newcommand{\JudShort}[1]{\J_{#1}}

% Components
\newcommand{\coord}{\mathsf{c}}            % coordinate
\newcommand{\prop}{\varphi}                 % proposition
\newcommand{\carrier}{\mathcal{A}}          % carrier type
\newcommand{\evid}{\mathcal{E}}             % evidence bundle
\newcommand{\oblig}{\mathcal{O}}            % obligations
\newcommand{\obstr}{\mathcal{B}}            % obstructions
\newcommand{\trust}{\mathcal{T}}            % trust annotation
\newcommand{\prov}{\Pi}                     % provenance

% Trust algebra
\newcommand{\Talg}{\mathfrak{T}}            % trust ordered algebra
\newcommand{\Eadm}{\mathcal{E}_{\mathrm{adm}}}
\newcommand{\tleq}{\preceq}                 % trust ordering
\newcommand{\tjoin}{\oplus}                 % conservative join
\newcommand{\tatten}{\ominus}               % attenuation
\newcommand{\tpromote}{\uparrow_\pi}        % promotion
\newcommand{\tdemote}{\downarrow_\chi}       % demotion

% Trust tiers
\newcommand{\Tcontra}{\ensuremath{\mathsf{CONTRADICTED}}}
\newcommand{\Tunver}{\ensuremath{\mathsf{UNVERIFIED}}}
\newcommand{\Tcopilot}{\ensuremath{\mathsf{COPILOT}}}
\newcommand{\Toracle}{\ensuremath{\mathsf{ORACLE}}}
\newcommand{\Truntime}{\ensuremath{\mathsf{RUNTIME}}}
\newcommand{\Tsolver}{\ensuremath{\mathsf{SOLVER}}}
\newcommand{\Tproof}{\ensuremath{\mathsf{PROOF}}}

% Site and geometry
\newcommand{\Site}{\mathcal{S}}             % semantic site
\newcommand{\Cov}{\mathrm{Cov}}            % covering family
\newcommand{\Ob}{\mathrm{Ob}}              % objects
\newcommand{\Mor}{\mathrm{Mor}}            % morphisms
\newcommand{\res}{\mathrm{res}}            % restriction
\newcommand{\Cech}{\check{C}}              % Čech complex
\newcommand{\Hcoh}[1]{H^{#1}}             % cohomology group

% Descent
\newcommand{\descent}{\mathrm{desc}}
\newcommand{\glue}{\mathrm{glue}}
\newcommand{\overlap}[2]{#1 \cap #2}

% Semantic moves
\newcommand{\VERIFY}{\textsc{Verify}}
\newcommand{\CONSTRUCT}{\textsc{Construct}}
\newcommand{\NORMALIZE}{\textsc{Normalize}}
\newcommand{\GLUE}{\textsc{Glue}}
\newcommand{\RESTRICT}{\textsc{Restrict}}
\newcommand{\TRANSPORT}{\textsc{Transport}}
\newcommand{\REPAIR}{\textsc{Repair}}
\newcommand{\PROMOTE}{\textsc{Promote}}
\newcommand{\CHALLENGE}{\textsc{Challenge}}
\newcommand{\ESCALATE}{\textsc{Escalate}}

% Fragments
\newcommand{\QFLIA}{\texttt{QF\_LIA}}
\newcommand{\QFLRA}{\texttt{QF\_LRA}}
\newcommand{\QFBV}{\texttt{QF\_BV}}
\newcommand{\QFUF}{\texttt{QF\_UF}}

% Misc
\newcommand{\Python}{\textsc{Python}}
\newcommand{\JuGeo}{\textsc{JuGeo}}
\newcommand{\LEAN}{\textsc{Lean}\,4}
\newcommand{\Fstar}{F$^{\star}$}
\newcommand{\Coq}{\textsc{Coq}}
\newcommand{\Dafny}{\textsc{Dafny}}

% Common shortcuts
\newcommand{\ie}{i.e.\@\xspace}
\newcommand{\eg}{e.g.\@\xspace}
\newcommand{\cf}{cf.\@\xspace}
\newcommand{\wrt}{w.r.t.\@\xspace}

% Evaluation shortcuts
\newcommand{\numcases}{300}
\newcommand{\numspec}{100}
\newcommand{\numeq}{100}
\newcommand{\numbug}{100}
\newcommand{\medtime}{6.5\,ms}
\newcommand{\totaltime}{1.949\,s}


% ── Packages ────────────────────────────────────────────────────────────
\usepackage{amsmath,amssymb,amsthm,mathtools}
\usepackage{tikz-cd}
\usepackage{listings}
\usepackage{xcolor}
\usepackage{xspace}
\usepackage{booktabs}
\usepackage{multirow}
\usepackage{enumitem}
\usepackage[expansion=false]{microtype}
\usepackage{hyperref}
\usepackage{cleveref}
% \usepackage{thmtools}  % disabled: not available in this TeX installation
\usepackage{float}
\usepackage{caption}
\usepackage{subcaption}
\usepackage{tabularx}
\usepackage{stmaryrd}       % \llbracket, \rrbracket
\usepackage{bbm}            % \mathbbm{1}

% ── Page geometry ───────────────────────────────────────────────────────
\usepackage[margin=1in]{geometry}

% ── Theorem environments ────────────────────────────────────────────────
\theoremstyle{plain}
\newtheorem{theorem}{Theorem}[section]
\newtheorem{lemma}[theorem]{Lemma}
\newtheorem{proposition}[theorem]{Proposition}
\newtheorem{corollary}[theorem]{Corollary}

\theoremstyle{definition}
\newtheorem{definition}[theorem]{Definition}
\newtheorem{example}[theorem]{Example}
\newtheorem{construction}[theorem]{Construction}

\theoremstyle{remark}
\newtheorem{remark}[theorem]{Remark}
\newtheorem{notation}[theorem]{Notation}

% ── Python listings style ──────────────────────────────────────────────
\definecolor{jugeoBlue}{HTML}{2563EB}
\definecolor{jugeoGreen}{HTML}{059669}
\definecolor{jugeoRed}{HTML}{DC2626}
\definecolor{jugeoPurple}{HTML}{7C3AED}
\definecolor{jugeoGray}{HTML}{6B7280}
\definecolor{jugeoBg}{HTML}{F8FAFC}

\lstdefinestyle{jugeo-python}{
  language=Python,
  basicstyle=\ttfamily\small,
  keywordstyle=\color{jugeoBlue}\bfseries,
  stringstyle=\color{jugeoGreen},
  commentstyle=\color{jugeoGray}\itshape,
  numberstyle=\tiny\color{jugeoGray},
  backgroundcolor=\color{jugeoBg},
  numbers=left,
  numbersep=6pt,
  frame=single,
  framerule=0.4pt,
  rulecolor=\color{jugeoGray!40},
  breaklines=true,
  breakatwhitespace=true,
  showstringspaces=false,
  tabsize=4,
  xleftmargin=1.5em,
  framexleftmargin=1.5em,
  morekeywords={dataclass,frozen,field,Enum,IntEnum,auto,Any,
                Mapping,Sequence,Optional,match,case},
  literate={→}{{$\to$}}1 {←}{{$\leftarrow$}}1
           {≤}{{$\leq$}}1 {≥}{{$\geq$}}1 {≠}{{$\neq$}}1
           {∈}{{$\in$}}1 {∀}{{$\forall$}}1 {∃}{{$\exists$}}1
           {⊕}{{$\oplus$}}1 {⊖}{{$\ominus$}}1,
  escapeinside={(*@}{@*)},
}
\lstset{style=jugeo-python}

\lstdefinestyle{jugeo-output}{
  basicstyle=\ttfamily\small,
  backgroundcolor=\color{jugeoBg},
  frame=single,
  framerule=0.4pt,
  rulecolor=\color{jugeoGray!40},
  breaklines=true,
  showstringspaces=false,
  xleftmargin=1.5em,
  framexleftmargin=0.5em,
  numbers=none,
}

% ── JuGeo-specific macros ──────────────────────────────────────────────

% Judgment 8-tuple
\newcommand{\J}{\mathcal{J}}
\newcommand{\Jud}[8]{(#1,\,#2,\,#3,\,#4,\,#5,\,#6,\,#7,\,#8)}
\newcommand{\JudShort}[1]{\J_{#1}}

% Components
\newcommand{\coord}{\mathsf{c}}            % coordinate
\newcommand{\prop}{\varphi}                 % proposition
\newcommand{\carrier}{\mathcal{A}}          % carrier type
\newcommand{\evid}{\mathcal{E}}             % evidence bundle
\newcommand{\oblig}{\mathcal{O}}            % obligations
\newcommand{\obstr}{\mathcal{B}}            % obstructions
\newcommand{\trust}{\mathcal{T}}            % trust annotation
\newcommand{\prov}{\Pi}                     % provenance

% Trust algebra
\newcommand{\Talg}{\mathfrak{T}}            % trust ordered algebra
\newcommand{\Eadm}{\mathcal{E}_{\mathrm{adm}}}
\newcommand{\tleq}{\preceq}                 % trust ordering
\newcommand{\tjoin}{\oplus}                 % conservative join
\newcommand{\tatten}{\ominus}               % attenuation
\newcommand{\tpromote}{\uparrow_\pi}        % promotion
\newcommand{\tdemote}{\downarrow_\chi}       % demotion

% Trust tiers
\newcommand{\Tcontra}{\ensuremath{\mathsf{CONTRADICTED}}}
\newcommand{\Tunver}{\ensuremath{\mathsf{UNVERIFIED}}}
\newcommand{\Tcopilot}{\ensuremath{\mathsf{COPILOT}}}
\newcommand{\Toracle}{\ensuremath{\mathsf{ORACLE}}}
\newcommand{\Truntime}{\ensuremath{\mathsf{RUNTIME}}}
\newcommand{\Tsolver}{\ensuremath{\mathsf{SOLVER}}}
\newcommand{\Tproof}{\ensuremath{\mathsf{PROOF}}}

% Site and geometry
\newcommand{\Site}{\mathcal{S}}             % semantic site
\newcommand{\Cov}{\mathrm{Cov}}            % covering family
\newcommand{\Ob}{\mathrm{Ob}}              % objects
\newcommand{\Mor}{\mathrm{Mor}}            % morphisms
\newcommand{\res}{\mathrm{res}}            % restriction
\newcommand{\Cech}{\check{C}}              % Čech complex
\newcommand{\Hcoh}[1]{H^{#1}}             % cohomology group

% Descent
\newcommand{\descent}{\mathrm{desc}}
\newcommand{\glue}{\mathrm{glue}}
\newcommand{\overlap}[2]{#1 \cap #2}

% Semantic moves
\newcommand{\VERIFY}{\textsc{Verify}}
\newcommand{\CONSTRUCT}{\textsc{Construct}}
\newcommand{\NORMALIZE}{\textsc{Normalize}}
\newcommand{\GLUE}{\textsc{Glue}}
\newcommand{\RESTRICT}{\textsc{Restrict}}
\newcommand{\TRANSPORT}{\textsc{Transport}}
\newcommand{\REPAIR}{\textsc{Repair}}
\newcommand{\PROMOTE}{\textsc{Promote}}
\newcommand{\CHALLENGE}{\textsc{Challenge}}
\newcommand{\ESCALATE}{\textsc{Escalate}}

% Fragments
\newcommand{\QFLIA}{\texttt{QF\_LIA}}
\newcommand{\QFLRA}{\texttt{QF\_LRA}}
\newcommand{\QFBV}{\texttt{QF\_BV}}
\newcommand{\QFUF}{\texttt{QF\_UF}}

% Misc
\newcommand{\Python}{\textsc{Python}}
\newcommand{\JuGeo}{\textsc{JuGeo}}
\newcommand{\LEAN}{\textsc{Lean}\,4}
\newcommand{\Fstar}{F$^{\star}$}
\newcommand{\Coq}{\textsc{Coq}}
\newcommand{\Dafny}{\textsc{Dafny}}

% Common shortcuts
\newcommand{\ie}{i.e.\@\xspace}
\newcommand{\eg}{e.g.\@\xspace}
\newcommand{\cf}{cf.\@\xspace}
\newcommand{\wrt}{w.r.t.\@\xspace}

% Evaluation shortcuts
\newcommand{\numcases}{300}
\newcommand{\numspec}{100}
\newcommand{\numeq}{100}
\newcommand{\numbug}{100}
\newcommand{\medtime}{6.5\,ms}
\newcommand{\totaltime}{1.949\,s}


% ——— Paper-specific macros (letters only) ———
\newcommand{\CDA}{\mathsf{CopilotDeductionAssist}}
\newcommand{\DedHint}{\mathrm{deductionHint}}
\newcommand{\DedPlan}{\mathrm{deductionPlan}}
\newcommand{\DedStep}{\mathrm{deductionStep}}
\newcommand{\DedTrust}{\mathrm{hintTrust}}
\newcommand{\DedCalib}{\mathrm{calibration}}
\newcommand{\RuleApp}{\mathrm{ruleApplication}}
\newcommand{\SubGoal}{\mathrm{subGoal}}
\newcommand{\CopChan}{\mathsf{CopilotChannel}}
\newcommand{\CopBridge}{\mathsf{CopilotAPIBridge}}
\newcommand{\ChanPool}{\mathsf{ChannelPool}}
\newcommand{\ChanMon}{\mathsf{ChannelMonitor}}
\newcommand{\ChanRouter}{\mathsf{ChannelRouter}}
\newcommand{\EvReq}{\mathsf{EvidenceRequest}}
\newcommand{\EvResp}{\mathsf{EvidenceResponse}}

\title{\textbf{CopilotDeductionAssist:\\
Trust-Calibrated Deduction Hints in Judgment Geometry}}
\author{Halley Young}
\date{Paper~75 --- Judgment Geometry Series}

\begin{document}
\maketitle

% ============================================================
%  Abstract
% ============================================================
\begin{abstract}
Deduction---the application of inference rules to derive new
judgments from existing ones---is the fundamental mechanism by
which the \JuGeo{} kernel advances proofs.  While the kernel's
deduction engine is sound by construction, choosing \emph{which}
rules to apply, in what order, and with what arguments is a
search problem.  This paper introduces the $\CDA$, a Copilot-guided
assistant that provides trust-calibrated deduction hints.  The
assistant suggests rule applications, ranks deduction plans by
estimated utility, and provides calibrated confidence scores.
All hints operate at the $\Tcopilot$ trust tier and require
validation before installation.  We formalise deduction hints
as morphisms in a rule-application category, prove four main
theorems (hint soundness, trust calibration, deduction
completeness, and plan optimality), and provide four \Python{}
code listings.  An appendix contains additional lemmas on
category axioms and calibration properties.
\end{abstract}

% ============================================================
%  Section 1 — Introduction
% ============================================================
\section{Introduction}\label{sec:intro}

The deduction engine in \JuGeo{} applies inference rules
drawn from a rule base.  Given a goal, the engine searches
for a sequence of rule applications that derives the goal
from available hypotheses and previously proven lemmas.
This search can be expensive: the rule base may contain
hundreds of rules, each with multiple instantiation patterns.

The $\CDA$ addresses this by providing \emph{hints}:
ranked suggestions for which rules to try, in what order,
and with what arguments.  The assistant does not modify
the proof state; it merely guides the search.  The trust
model ensures that hints are treated as suggestions, not
as proofs.

\paragraph{Contributions.}
\begin{enumerate}
\item A categorical model of deduction as a rule-application
      category (\S\ref{sec:deduction}).
\item The design of the $\CDA$ with three capabilities:
      rule suggestion, plan ranking, and confidence
      calibration (\S\ref{sec:assistant}).
\item Four theorems establishing soundness, calibration,
      completeness, and optimality (\S\ref{sec:formal}).
\item Four \Python{} listings showing integration with
      the channel layer (\S\ref{sec:code}).
\end{enumerate}

\paragraph{Paper outline.}
Section~\ref{sec:background} gives background.
Section~\ref{sec:deduction} formalises deduction.
Section~\ref{sec:assistant} presents the assistant design.
Section~\ref{sec:code} shows the implementation.
Section~\ref{sec:formal} states the main theorems.
Section~\ref{sec:discussion} discusses limitations.
Section~\ref{sec:related} surveys related work.
Section~\ref{sec:conclusion} concludes.
Appendix~\ref{app:proofs} contains auxiliary proofs.

% ============================================================
%  Section 2 — Background
% ============================================================
\section{Background}\label{sec:background}

\subsection{Trust lattice}

The trust ordering:
\[
  \Tcontra \;\tleq\; \Tunver \;\tleq\; \Tcopilot
  \;\tleq\; \Toracle \;\tleq\; \Truntime
  \;\tleq\; \Tsolver \;\tleq\; \Tproof.
\]
The deduction assistant operates at $\Tcopilot$.

\subsection{Inference rules}

An \emph{inference rule} is a pair
$(\Gamma, \varphi)$ where $\Gamma = \{\psi_1, \ldots, \psi_n\}$
are the premises and $\varphi$ is the conclusion.  A rule is
\emph{applicable} to a goal $\varphi$ if $\varphi$ matches
the conclusion and all premises are derivable or are open goals.

\begin{definition}[Rule base]\label{def:rulebase}
A \emph{rule base} $\mathcal{R}$ is a finite set of inference
rules.  The kernel's deduction engine applies rules from
$\mathcal{R}$ to advance proofs.
\end{definition}

\begin{definition}[Rule application]\label{def:ruleapp}
A \emph{rule application} is a tuple $(r, \sigma, g)$ where
$r \in \mathcal{R}$ is a rule, $\sigma$ is a substitution
instantiating the rule's schematic variables, and $g$ is the
goal to which the rule is applied.
\end{definition}

\subsection{Deduction search}

\begin{definition}[Deduction tree]\label{def:dedtree}
A \emph{deduction tree} is a tree whose root is the goal,
internal nodes are rule applications, and leaves are either
axioms (closed) or open sub-goals.
\end{definition}

\begin{definition}[Deduction plan]\label{def:dedplan}
A \emph{deduction plan} is a sequence of rule applications
$(a_1, \ldots, a_k)$ that, when applied in order, transforms
the initial goal into a set of sub-goals (possibly empty).
A plan with no remaining sub-goals constitutes a complete
deduction.
\end{definition}

% ============================================================
%  Section 3 — Formalisation
% ============================================================
\section{Deduction Formalisation}\label{sec:deduction}

\subsection{Rule-application category}

\begin{definition}[Rule-application category]\label{def:racat}
Let $\mathbf{RA}$ be the category whose objects are proof
states (sets of open goals) and whose morphisms are rule
applications: $a \colon S \to S'$ means applying rule
application~$a$ to state~$S$ produces state~$S'$.
\end{definition}

\begin{lemma}[Category axioms]\label{lem:cataxioms}
$\mathbf{RA}$ satisfies the category axioms:
\begin{enumerate}
\item \emph{Identity:} the no-op application $\mathrm{id}_S$
      leaves state~$S$ unchanged.
\item \emph{Composition:} if $a_1 \colon S \to S'$ and
      $a_2 \colon S' \to S''$, then
      $a_2 \circ a_1 \colon S \to S''$.
\item \emph{Associativity:} $(a_3 \circ a_2) \circ a_1 =
      a_3 \circ (a_2 \circ a_1)$.
\end{enumerate}
\end{lemma}
\begin{proof}
\emph{(1)}~The no-op does not change the goal set.
\emph{(2)}~Applying $a_1$ then $a_2$ is well-defined since
$a_2$ applies to the state produced by $a_1$.
\emph{(3)}~Application is sequential; the result depends only
on the intermediate states, which are uniquely determined.
\end{proof}

\subsection{Hint as a morphism selector}

\begin{definition}[Deduction hint]\label{def:dedhint}
A \emph{deduction hint} is a function
$h \colon \mathsf{Obj}(\mathbf{RA}) \to
\mathcal{P}(\mathsf{Mor}(\mathbf{RA}))$
that assigns to each proof state a ranked set of candidate
rule applications.  Each candidate carries a confidence
score $c \in [0,1]$.
\end{definition}

\begin{definition}[Hint quality]\label{def:hintquality}
The \emph{quality} of a hint~$h$ at state~$S$ is:
\[
  Q(h, S) = \sum_{a \in h(S)} c(a) \cdot \mathbb{1}[\text{$a$
  leads to a shorter proof}].
\]
Higher quality means more correct, highly-ranked suggestions.
\end{definition}

\subsection{Calibration}

\begin{definition}[Calibrated hint]\label{def:calibhint}
A hint~$h$ is \emph{calibrated} if for every confidence
level~$c$, the fraction of suggestions with confidence
$\geq c$ that are correct is at least~$c$.  Formally:
\[
  \Pr[\text{suggestion correct} \mid c(\text{suggestion}) \geq c]
  \geq c.
\]
\end{definition}

\begin{remark}
Calibration ensures that the confidence scores are
meaningful: a suggestion with 90\% confidence should be
correct at least 90\% of the time.
\end{remark}

% ============================================================
%  Section 4 — The Deduction Assistant
% ============================================================
\section{The Deduction Assistant}\label{sec:assistant}

\subsection{Architecture}

The $\CDA$ has three capabilities:
\begin{enumerate}
\item \emph{Rule suggestion:} given a goal, suggest which
      inference rules to apply.
\item \emph{Plan ranking:} given multiple deduction plans,
      rank them by estimated quality.
\item \emph{Confidence calibration:} provide calibrated
      confidence scores for each suggestion.
\end{enumerate}

\subsection{Rule suggestion}

The assistant queries the Copilot channel with a prompt
encoding the goal, the available rules, and the proof context.
The response is a ranked list of rule applications.

\begin{definition}[Rule suggestion protocol]\label{def:rulesuggprot}
The rule suggestion protocol:
\begin{enumerate}
\item Encode the goal as a string.
\item Append the names and schemas of available rules.
\item Query the Copilot channel.
\item Parse the response into rule applications.
\item Filter invalid applications (those that do not match
      the goal).
\item Return the filtered list with confidence scores.
\end{enumerate}
\end{definition}

\subsection{Plan ranking}

The assistant ranks deduction plans by querying the Copilot
channel for an assessment of each plan's likely success.

\begin{definition}[Plan ranking function]\label{def:planrank}
The plan ranking function
$R \colon \mathcal{P}(\mathsf{Plan}) \to \mathsf{List}(\mathsf{Plan})$
orders plans by estimated quality.  The quality estimate
is the product of the confidence scores of the plan's steps.
\end{definition}

\subsection{Confidence calibration}

The assistant's confidence scores are calibrated using
historical data from the $\ChanMon$.  The calibration
function adjusts raw LLM confidence scores based on the
observed success rate at each confidence level.

\begin{definition}[Calibration function]\label{def:calibfunc}
The calibration function
$\kappa \colon [0,1] \to [0,1]$
maps raw confidence $c_{\mathrm{raw}}$ to calibrated
confidence $c_{\mathrm{cal}}$ using isotonic regression
on historical success data.
\end{definition}

\begin{lemma}[Calibration is monotone]\label{lem:calibmono}
$c_1 \leq c_2 \implies \kappa(c_1) \leq \kappa(c_2)$.
\end{lemma}
\begin{proof}
Isotonic regression produces a monotone function by
construction: it fits a non-decreasing step function to
the data.
\end{proof}

% ============================================================
%  Section 5 — Implementation
% ============================================================
\section{Implementation}\label{sec:code}

\subsection{Rule suggestion via CopilotChannel}

\begin{lstlisting}[style=jugeo-python,
  caption={Deduction hint generation via CopilotChannel},
  label={lst:dedhint}]
from jugeo.evidence.channels import (
    CopilotChannel,
    ChannelConfiguration,
    EvidenceChannel,
    EvidenceRequest,
    EvidenceResponse,
    ChannelJurisdiction,
    ChannelMonitor,
)

class CopilotDeductionAssist:
    """Copilot-guided deduction hint provider."""

    def __init__(self, channel: CopilotChannel, monitor: ChannelMonitor):
        self.channel = channel
        self.monitor = monitor
        self.trust_ceiling = channel.TRUST_CEILING
        self.jurisdiction = ChannelJurisdiction.for_channel(
            EvidenceChannel.COPILOT
        )
        self.calibration_data = {}

    def suggest_rules(self, goal, available_rules, context=None):
        prompt = self._build_rule_prompt(goal, available_rules, context)
        raw = self.channel.query_llm(prompt, timeout_ms=10000)
        parsed = self.channel.parse_response()
        clamped = self.channel.apply_trust_ceiling()
        suggestions = self._filter_and_rank(clamped, available_rules)
        calibrated = self._calibrate(suggestions)
        return calibrated

    def rank_plans(self, plans, goal):
        prompt = self._build_ranking_prompt(plans, goal)
        raw = self.channel.query_llm(prompt, timeout_ms=8000)
        parsed = self.channel.parse_response()
        clamped = self.channel.apply_trust_ceiling()
        return self._parse_ranking(clamped, plans)

    def update_calibration(self, suggestion, was_correct):
        conf = suggestion.get("confidence", 0.5)
        bucket = round(conf, 1)
        if bucket not in self.calibration_data:
            self.calibration_data[bucket] = {"correct": 0, "total": 0}
        self.calibration_data[bucket]["total"] += 1
        if was_correct:
            self.calibration_data[bucket]["correct"] += 1

    def _build_rule_prompt(self, goal, rules, ctx):
        return (
            f"Goal: {goal}. "
            f"Available rules: {[r['name'] for r in rules[:15]]}. "
            f"Context: {ctx}. "
            f"Suggest rule applications with confidence scores."
        )

    def _build_ranking_prompt(self, plans, goal):
        return f"Rank these deduction plans for goal {goal}: {plans}"

    def _filter_and_rank(self, response, rules):
        rule_names = {r["name"] for r in rules}
        if isinstance(response, list):
            return [s for s in response
                    if s.get("rule") in rule_names]
        return []

    def _calibrate(self, suggestions):
        for s in suggestions:
            raw_conf = s.get("confidence", 0.5)
            bucket = round(raw_conf, 1)
            data = self.calibration_data.get(bucket)
            if data and data["total"] >= 10:
                s["confidence"] = data["correct"] / data["total"]
        return suggestions

    def _parse_ranking(self, response, plans):
        return sorted(plans, key=lambda p: p.get("score", 0), reverse=True)
\end{lstlisting}

\subsection{Deduction with CopilotAPIBridge}

\begin{lstlisting}[style=jugeo-python,
  caption={Deduction orchestration with CopilotAPIBridge},
  label={lst:dedbridge}]
from jugeo.interfaces.api import (
    CopilotAPIBridge,
    APIEventLog,
    APIRequest,
    APIResponse,
    OperationKind,
    JuGeoAPI,
    APISession,
)
from jugeo.evidence.channels import (
    CopilotChannel,
    ChannelConfiguration,
    EvidenceChannel,
)

class DeductionOrchestrator:
    """Orchestrate deduction with Copilot assistance."""

    def __init__(self, bridge: CopilotAPIBridge, api: JuGeoAPI):
        self.bridge = bridge
        self.api = api
        self.event_log = APIEventLog()

    def deduce_with_hints(self, coordinate, proposition, rules):
        session = self.api.open_session("deduction")

        hint = self.bridge.query(
            f"Suggest deduction strategy for {proposition}"
        )
        self.bridge.enforce_trust_ceiling()
        self.event_log.record({
            "event": "deduction_hint",
            "hint": str(hint),
        })

        request = APIRequest(
            request_id="ded-001",
            operation=OperationKind.VERIFY,
            coordinate=coordinate,
            proposition=proposition,
            budget=30000,
            deadline=45000,
            metadata={"deduction_hint": str(hint), "rules": len(rules)},
        )
        session.register_request(request)
        result = self.api.verify(
            coordinate=coordinate,
            proposition=proposition,
        )

        self.event_log.record({
            "event": "deduction_result",
            "success": result.is_successful() if result else False,
        })
        session.checkpoint()

        if result and result.has_residuals():
            sub_hints = self.bridge.query(
                f"Handle sub-goals from deduction: {result}"
            )
            self.bridge.enforce_trust_ceiling()
            stats = self.bridge.stats()
            session.close()
            return {"result": result, "sub_hints": str(sub_hints),
                    "stats": stats}

        session.close()
        return {"result": result}
\end{lstlisting}

\subsection{Deduction evidence request flow}

\begin{lstlisting}[style=jugeo-python,
  caption={Evidence request flow for deduction hints},
  label={lst:dedevflow}]
from jugeo.evidence.channels import (
    EvidenceRequest,
    EvidenceResponse,
    EvidenceRecord,
    SupportRoute,
    EvidenceChannel,
    CopilotChannel,
    SolverChannel,
    ChannelRouter,
    ChannelPool,
    ChannelMonitor,
)

def deduction_evidence_flow(
    pool: ChannelPool,
    router: ChannelRouter,
    goal,
    rule_application,
):
    """Full evidence flow for a deduction hint validation."""
    monitor = ChannelMonitor()

    request = EvidenceRequest(
        request_id=f"ded-ev-{id(goal)}",
        coordinate=str(goal),
        proposition=f"rule {rule_application} applies to {goal}",
        required_kind=None,
        proposition_kind="deduction_validation",
        preferred_channel=EvidenceChannel.SOLVER,
        fallback_channels=[EvidenceChannel.COPILOT],
        deadline_ms=15000,
        budget=20000,
        metadata={"rule": str(rule_application)},
    )

    if request.is_expired():
        return None
    remaining = request.remaining_ms()
    canonical = request.canonical_key()

    best = router.find_best_channel(request)
    all_capable = router.find_all_capable(request)

    monitor.record_request(request)

    solver = pool.get_channel("solver")
    if solver:
        solver_resp = solver.handle_request(request)
        monitor.record_response(solver_resp)
        if solver_resp and not solver_resp.is_empty():
            record = EvidenceRecord(
                record_id=f"ded-rec-{id(goal)}",
                request=request,
                response=solver_resp,
                channel=EvidenceChannel.SOLVER,
                timestamp=None,
            )
            route = SupportRoute(
                source=EvidenceChannel.SOLVER,
                target=str(goal),
                trust="solver",
            )
            record_with_route = record.with_support_route(route)
            return {
                "record": record_with_route,
                "regions": record_with_route.support_regions(),
                "canonical": record_with_route.canonical_key(),
                "success_rate": monitor.success_rate(),
            }

    copilot = pool.get_channel("copilot")
    if copilot:
        copilot_resp = copilot.handle_request(request)
        monitor.record_response(copilot_resp)
        return {"fallback": True, "response": copilot_resp}

    return None
\end{lstlisting}

\subsection{Deduction monitoring and diagnostics}

\begin{lstlisting}[style=jugeo-python,
  caption={Deduction monitoring and diagnostics},
  label={lst:deddiag}]
from jugeo.kernel.health import (
    HealthStatus,
    HealthDimension,
    HealthIndicator,
    HealthReport,
)
from jugeo.evidence.channels import (
    ChannelPool,
    ChannelMonitor,
    CopilotChannel,
)
from jugeo.interfaces.api import (
    JuGeoAPI,
    APIEventLog,
)

class DeductionDiagnostics:
    """Monitor deduction assistant performance."""

    def __init__(self, pool: ChannelPool, monitor: ChannelMonitor):
        self.pool = pool
        self.monitor = monitor

    def hint_success_rate(self):
        return self.monitor.success_rate()

    def hint_latency(self):
        return self.monitor.latency_percentiles()

    def copilot_health(self):
        copilot = self.pool.get_channel("copilot")
        if copilot is None:
            return HealthIndicator(
                dimension=HealthDimension.COPILOT_CONNECTIVITY,
                status=HealthStatus.UNHEALTHY,
                message="No copilot channel",
                value=0.0,
                timestamp=None,
                details={},
            )
        alive = copilot.health_check()
        return HealthIndicator(
            dimension=HealthDimension.COPILOT_CONNECTIVITY,
            status=HealthStatus.HEALTHY if alive
                   else HealthStatus.DEGRADED,
            message="Copilot for deduction",
            value=1.0 if alive else 0.5,
            timestamp=None,
            details={"alive": alive},
        )

    def full_report(self):
        copilot_ind = self.copilot_health()
        success = self.hint_success_rate()
        flow_ind = HealthIndicator(
            dimension=HealthDimension.EVIDENCE_FLOW,
            status=(HealthStatus.HEALTHY if success >= 0.8
                    else HealthStatus.DEGRADED),
            message=f"Deduction hint success: {success:.2%}",
            value=success,
            timestamp=None,
            details=self.monitor.summary(),
        )
        indicators = [copilot_ind, flow_ind]
        overall = (HealthStatus.HEALTHY
                   if all(i.is_healthy() for i in indicators)
                   else HealthStatus.DEGRADED)
        return HealthReport(
            overall_status=overall,
            indicators=indicators,
            generated_at=None,
            report_id="ded-health",
            degradation_reasons=[],
            recovery_suggestions=[],
        )
\end{lstlisting}

% ============================================================
%  Section 6 — Formal Properties
% ============================================================
\section{Formal Properties}\label{sec:formal}

\begin{theorem}[Hint soundness]\label{thm:hintsound}
Every deduction hint produced by the $\CDA$ is either valid
(the suggested rule application is applicable) or is filtered
out during the validation step.  No invalid hint reaches the
deduction engine.
\end{theorem}
\begin{proof}
The assistant's rule suggestion protocol
(Definition~\ref{def:rulesuggprot}) includes a filtering step
that checks each suggested rule application against the current
goal.  Applications that do not match are discarded.  Those
that pass the filter are submitted to the solver channel for
validation.  Invalid applications are rejected by the solver.
Therefore, only valid applications reach the engine.
\end{proof}

\begin{theorem}[Trust calibration]\label{thm:trustcalib}
The $\CDA$ produces calibrated confidence scores (in the limit).
Specifically, after sufficiently many interactions, the
calibration function $\kappa$ satisfies
Definition~\ref{def:calibhint}.
\end{theorem}
\begin{proof}
The calibration function uses isotonic regression on historical
data.  By the law of large numbers, as the number of suggestions
grows, the observed success rate at each confidence bucket
converges to the true probability.  Isotonic regression fits
a monotone function to these converging estimates.  In the limit,
the calibration satisfies Definition~\ref{def:calibhint}.
\end{proof}

\begin{theorem}[Deduction completeness]\label{thm:dedcomp}
If a goal is derivable from the rule base, then the $\CDA$
eventually suggests a complete deduction plan (in the idealised
setting where the language model is complete).
\end{theorem}
\begin{proof}
If the goal is derivable, there exists a finite deduction plan
$\pi = (a_1, \ldots, a_k)$.  In the idealised setting, the
language model enumerates all rule applications for each sub-goal,
including those in~$\pi$.  The plan ranking function assigns
positive quality to~$\pi$ (since all its steps are correct).
Therefore~$\pi$ is included in the ranked plan list.
\end{proof}

\begin{theorem}[Plan optimality]\label{thm:planopt}
Among the plans suggested by the $\CDA$, the highest-ranked
plan has the highest estimated quality (by
Definition~\ref{def:hintquality}).
\end{theorem}
\begin{proof}
The plan ranking function sorts plans by estimated quality
in descending order.  The first plan in the sorted list has
the highest estimated quality.  This is optimal with respect
to the estimate; the actual quality may differ, but the ranking
is correct with respect to the available information.
\end{proof}

% ============================================================
%  Section 7 — Discussion
% ============================================================
\section{Discussion}\label{sec:discussion}

\paragraph{Calibration warm-up.}
The calibration function requires historical data to be
accurate.  During the warm-up period (first ~100 suggestions),
the raw confidence scores from the LLM are used without
calibration.  The assistant flags these as uncalibrated.

\paragraph{Rule base size.}
Large rule bases (hundreds of rules) make rule suggestion
challenging.  The assistant mitigates this by filtering rules
by type match before querying the Copilot channel, reducing
the search space.

\paragraph{Interaction with construction participant.}
The deduction assistant and construction participant (Paper~74)
are complementary: the construction participant generates proof
plans, while the deduction assistant suggests individual rule
applications within those plans.

\paragraph{Deduction depth.}
For deep deductions (many rule applications), the assistant's
confidence in individual steps may compound unfavourably.
The plan ranking function accounts for this by penalising
longer plans.

\paragraph{Fallback to exhaustive search.}
When the assistant's suggestions fail, the deduction engine
falls back to exhaustive search over the rule base.  The
assistant's hints reduce the expected search time but do not
eliminate the fallback.

% ============================================================
%  Section 8 — Related Work
% ============================================================
\section{Related Work}\label{sec:related}

\paragraph{Tactic suggestion.}
Proverbot9001~\cite{sanchez2020proverbot} and
ASTactic~\cite{yang2019learning} suggest proof tactics using
machine learning.  Our assistant differs in using a trust
lattice and channel architecture.

\paragraph{Proof hint systems.}
Isabelle's \texttt{try} command~\cite{wenzel2014isabelle}
and Coq's \texttt{auto}~\cite{bertot2004coq} provide
automated deduction hints.  Our approach externalises hint
generation to the Copilot channel.

\paragraph{Confidence calibration.}
Platt scaling~\cite{platt1999probabilities} and isotonic
regression~\cite{zadrozny2002transforming} are standard
calibration methods.  We adopt isotonic regression for its
monotonicity guarantee.

\paragraph{Copilot advisors.}
Papers~74 and~80--84 of this series introduce related
advisors.  The deduction assistant is the most fundamental:
all other advisors ultimately rely on deduction for
validation.

% ============================================================
%  Section 9 — Conclusion
% ============================================================
\section{Conclusion}\label{sec:conclusion}

We have introduced the $\CDA$, a Copilot-guided assistant
that provides trust-calibrated deduction hints in the
\JuGeo{} framework.  The assistant suggests rule applications,
ranks deduction plans, and provides calibrated confidence
scores, all at $\Tcopilot$ trust.  Four theorems establish
hint soundness, trust calibration, deduction completeness,
and plan optimality.

Future work includes integrating the calibration function
with the channel monitor for online calibration updates,
supporting multi-step deduction hints (suggesting sequences
rather than individual rules), and adapting the assistant
to different logics.

% ============================================================
%  Bibliography
% ============================================================
\begin{thebibliography}{25}

\bibitem{sanchez2020proverbot}
A.~Sanchez-Stern, Y.~Alhessi, L.~Saul, and S.~Lerner,
``Generating correctness proofs with neural networks,''
\emph{MAPL}, 2020.

\bibitem{yang2019learning}
K.~Yang and J.~Deng,
``Learning to prove theorems via interacting with proof
assistants,''
\emph{ICML}, 2019.

\bibitem{wenzel2014isabelle}
M.~Wenzel,
``Isabelle/jEdit --- a prover IDE within the PIDE framework,''
\emph{AISC/Calculemus/MKM}, 2012.

\bibitem{bertot2004coq}
Y.~Bertot and P.~Cast\'{e}ran,
\emph{Interactive Theorem Proving and Program Development:
Coq'Art}.
Springer, 2004.

\bibitem{platt1999probabilities}
J.~Platt,
``Probabilistic outputs for support vector machines and
comparisons to regularized likelihood methods,''
\emph{Advances in Large Margin Classifiers}, 1999.

\bibitem{zadrozny2002transforming}
B.~Zadrozny and C.~Elkan,
``Transforming classifier scores into accurate multiclass
probability estimates,''
\emph{KDD}, 2002.

\bibitem{polu2020gptf}
S.~Polu and I.~Sutskever,
``Generative language modeling for automated theorem proving,''
\emph{arXiv preprint arXiv:2009.03393}, 2020.

\bibitem{lample2022hypertree}
G.~Lample \etal{},
``HyperTree proof search for neural theorem proving,''
\emph{NeurIPS}, 2022.

\bibitem{young2024jugeo}
H.~Young,
``Judgment geometry: a framework for trust-calibrated
verification,''
\emph{JuGeo Technical Reports}, Paper~1, 2024.

\bibitem{young2024channels}
H.~Young,
``Evidence channels in judgment geometry,''
\emph{JuGeo Technical Reports}, Paper~88, 2024.

\bibitem{young2024construction}
H.~Young,
``CopilotConstructionParticipant --- LLM-guided proof
construction,''
\emph{JuGeo Technical Reports}, Paper~74, 2024.

\bibitem{young2024advisors}
H.~Young,
``Copilot advisors for heap, scope, import, and callable
analysis,''
\emph{JuGeo Technical Reports}, Papers~80--81, 2024.

\bibitem{young2024lifecycle}
H.~Young,
``Kernel lifecycle and health monitoring in judgment geometry,''
\emph{JuGeo Technical Reports}, Paper~93, 2024.

\bibitem{demoura2008z3}
L.~de~Moura and N.~Bj{\o}rner,
``Z3: an efficient SMT solver,''
\emph{TACAS}, 2008.

\bibitem{moura2021lean4}
L.~de~Moura and S.~Ullrich,
``The Lean~4 theorem prover and programming language,''
\emph{CADE}, 2021.

\bibitem{bansal2019holist}
K.~Bansal, S.~Loos, M.~Rabe, C.~Szegedy, and S.~Wilcox,
``HOList: an environment for machine learning of higher-order
logic theorem proving,''
\emph{ICML}, 2019.

\bibitem{first2022diversity}
E.~First, M.~Rabe, T.~Ringer, and Y.~Brun,
``Diversity-driven automated formal verification,''
\emph{ICSE}, 2022.

\bibitem{jiang2022thor}
A.~Q.~Jiang, W.~Li, S.~Tworber, and Y.~Wu,
``Thor: wielding hammers to integrate language models and
automated theorem provers,''
\emph{NeurIPS}, 2022.

\bibitem{rabe2021mathematical}
M.~N.~Rabe, D.~Lee, K.~Bansal, and C.~Szegedy,
``Mathematical reasoning via self-supervised skip-tree
training,''
\emph{ICLR}, 2021.

\bibitem{li2021isarstep}
W.~Li, L.~Yu, Y.~Wu, and L.~C.~Paulson,
``IsarStep: a benchmark for high-level mathematical
reasoning,''
\emph{ICLR}, 2021.

\end{thebibliography}

% ============================================================
%  Appendix
% ============================================================
\appendix

\section{Auxiliary Proofs}\label{app:proofs}

\subsection{Category Axioms}

\begin{lemma}[Identity rule application is unique]%
\label{lem:idunique}
For each proof state~$S$, there is exactly one identity
morphism $\mathrm{id}_S$.
\end{lemma}
\begin{proof}
The identity morphism makes no changes to the goal set.
It is unique because the only function that leaves a set
unchanged is the identity function.
\end{proof}

\begin{lemma}[Composition is well-defined]%
\label{lem:compwd}
If $a_1 \colon S \to S'$ and $a_2 \colon S' \to S''$ are
rule applications, then $a_2 \circ a_1 \colon S \to S''$ is
a valid rule application sequence.
\end{lemma}
\begin{proof}
$a_1$ transforms $S$ to $S'$.  $a_2$ is defined on $S'$.
Their composition applies $a_1$ first, then $a_2$, producing
$S''$.  The composition is well-defined because the codomain
of $a_1$ equals the domain of $a_2$.
\end{proof}

\subsection{Calibration Properties}

\begin{lemma}[Isotonic regression preserves monotonicity]%
\label{lem:isomono}
The isotonic regression function $\kappa$ is non-decreasing.
\end{lemma}
\begin{proof}
By construction, isotonic regression fits a non-decreasing
step function to the data.  This is guaranteed by the
pool-adjacent-violators algorithm.
\end{proof}

\begin{lemma}[Calibration error converges to zero]%
\label{lem:caliberr}
As the number of observations $n \to \infty$, the expected
calibration error $\mathrm{ECE}$ converges to zero.
\end{lemma}
\begin{proof}
By the law of large numbers, the observed success rate in
each bucket converges to the true probability.  Isotonic
regression then maps each bucket to its true probability.
The expected calibration error, which measures the
discrepancy between predicted and observed success rates,
converges to zero.
\end{proof}

\begin{lemma}[Hint trust ceiling after calibration]%
\label{lem:hintceilcalib}
Calibration does not affect trust: the trust of a hint
remains $\tleq \Tcopilot$ regardless of the calibrated
confidence score.
\end{lemma}
\begin{proof}
Calibration adjusts the confidence score, not the trust
tier.  The trust is determined by the channel (Copilot)
and is clamped at $\Tcopilot$.  Confidence and trust are
independent properties.
\end{proof}

\subsection{Deduction Search Properties}

\begin{lemma}[Search space with hints is smaller]%
\label{lem:searchhint}
If the assistant suggests $k$ rule applications out of $n$
available, the search space at each step is reduced from $n$
to $k$ (where $k \leq n$).
\end{lemma}
\begin{proof}
Without hints, the engine considers all $n$ rules.  With
hints, it first tries the $k$ suggested rules (in priority
order), only falling back to the remaining $n - k$ rules
if none succeed.  In the best case, the effective search
space is $k$.
\end{proof}

\begin{lemma}[Hint-guided search terminates]%
\label{lem:hintsearchterm}
If the rule base is finite and the deduction depth is bounded,
then hint-guided search terminates.
\end{lemma}
\begin{proof}
At each step, the engine tries at most $n$ rules (with hints
reducing the expected number).  The depth is bounded by~$d$.
Therefore the total number of rule applications is at most
$n^d$, which is finite.
\end{proof}

\begin{lemma}[Plan quality monotonicity]%
\label{lem:planqualmono}
Adding a correct step to a plan does not decrease its quality:
$Q(\pi \cdot a, S) \geq Q(\pi, S)$ when $a$ is correct.
\end{lemma}
\begin{proof}
$Q$ is the sum of confidence-weighted correctness indicators.
A correct step contributes a positive term $c(a) > 0$.
Therefore the quality increases or stays the same.
\end{proof}

\begin{lemma}[Fallback completeness]%
\label{lem:fallbackcomp}
Even if all hints fail, the engine's exhaustive fallback
finds the derivation (if one exists), since the fallback
enumerates all rule applications up to the depth bound.
\end{lemma}
\begin{proof}
The fallback tries all $n$ rules at each depth level up
to~$d$.  If a derivation of depth $\leq d$ exists, the
exhaustive search finds it.
\end{proof}

\begin{proposition}[End-to-end deduction correctness]%
\label{prop:e2eded}
The deduction pipeline with Copilot hints produces only
correct derivations, regardless of the quality of the hints.
\end{proposition}
\begin{proof}
Hints guide the search but do not bypass validation.  Every
rule application is checked by the solver channel (at
$\Tsolver$).  Invalid applications are rejected.  Therefore,
only correct derivations are accepted.
\end{proof}

\begin{remark}[Multi-step hints]
The current assistant suggests one rule application at a time.
A multi-step hint would suggest a sequence of applications,
reducing the number of Copilot queries.  The trust model
extends naturally: each step in the sequence has trust
$\Tcopilot$, and the sequence trust is $\Tcopilot$.
\end{remark}

\section{Extended Formal Development}\label{app:extended}

\subsection{Rule-Application Algebra}

\begin{definition}[Rule-application algebra]%
\label{def:raalg}
A \emph{rule-application algebra} is a tuple
$(\mathcal{A}, \circ, e, \tau)$ where $\mathcal{A}$ is the set
of rule applications, $\circ$ is sequential composition,
$e$ is the identity (no-op), and $\tau$ assigns a trust tier.
\end{definition}

\begin{lemma}[Rule-application algebra is a monoid]%
\label{lem:raalgmonoid}
$(\mathcal{A}, \circ, e)$ is a monoid.
\end{lemma}
\begin{proof}
Associativity: sequential application of rules is function
composition, which is associative.
Identity: the no-op leaves the proof state unchanged.
\end{proof}

\begin{lemma}[Trust homomorphism for rule applications]%
\label{lem:trusthomra}
$\tau(a_1 \circ a_2) = \min(\tau(a_1), \tau(a_2))$.
\end{lemma}
\begin{proof}
Composed rule applications inherit the minimum trust of
their components, since the composition cannot produce
evidence above the lowest component's ceiling.
\end{proof}

\subsection{Deduction-Space Topology}

\begin{definition}[Deduction neighbourhood]%
\label{def:dedneighbour}
The \emph{deduction neighbourhood} of a proof state~$S$ is
the set of states reachable by a single rule application:
$N(S) = \{S' \mid \exists a.\; a \colon S \to S'\}$.
\end{definition}

\begin{lemma}[Neighbourhood is finite]%
\label{lem:neighfinite}
For a finite rule base $\mathcal{R}$ and a proof state~$S$
with finitely many open goals, $N(S)$ is finite.
\end{lemma}
\begin{proof}
Each rule in $\mathcal{R}$ can be applied to each open goal
with finitely many substitutions (assuming the substitution
space is finite, \eg{} for ground terms).  Therefore $|N(S)|
\leq |\mathcal{R}| \times |S| \times |\Sigma|$ where
$|\Sigma|$ is the number of valid substitutions.
\end{proof}

\begin{definition}[Deduction distance]%
\label{def:deddist}
The \emph{deduction distance} $d(S, S')$ between two proof
states is the minimum number of rule applications needed to
reach $S'$ from $S$.  If $S'$ is unreachable, $d(S, S') = \infty$.
\end{definition}

\begin{lemma}[Distance is a metric on reachable states]%
\label{lem:distmetric}
On the set of mutually reachable states, $d$ is a metric:
\begin{enumerate}
\item $d(S, S) = 0$,
\item $d(S, S') = d(S', S)$ (assuming reversible rules), and
\item $d(S, S'') \leq d(S, S') + d(S', S'')$ (triangle inequality).
\end{enumerate}
\end{lemma}
\begin{proof}
\emph{(1)}~Zero applications take $S$ to itself.
\emph{(2)}~If rules are reversible, any path from $S$ to $S'$
can be reversed.
\emph{(3)}~A path from $S$ to $S''$ via $S'$ has length
$d(S, S') + d(S', S'')$, which is at least the minimum
path length $d(S, S'')$.
\end{proof}

\subsection{Hint Effectiveness Measures}

\begin{definition}[Hint effectiveness]%
\label{def:hinteffect}
The \emph{effectiveness} of a hint $h$ at state $S$ is the
ratio of the search cost without hints to the search cost
with hints:
$E(h, S) = T_{\mathrm{no\text{-}hint}}(S) / T_{\mathrm{hint}}(S)$.
\end{definition}

\begin{lemma}[Effectiveness is at least 1]%
\label{lem:effectatleast}
$E(h, S) \geq 1$ for any hint $h$, since hints can only
reduce search cost (or leave it unchanged if ignored).
\end{lemma}
\begin{proof}
The hint-guided search tries the hinted rules first.  If
they fail, it falls back to exhaustive search.  In the worst
case, the hint adds a small constant overhead to the
exhaustive search.  Under the convention that we measure
asymptotic cost, $T_{\mathrm{hint}} \leq T_{\mathrm{no\text{-}hint}}$,
so $E \geq 1$.
\end{proof}

\begin{lemma}[Optimal hint effectiveness]%
\label{lem:opteffect}
If the hint always suggests the correct rule application
first, the effectiveness is $E(h, S) = n$, where $n = |\mathcal{R}|$.
\end{lemma}
\begin{proof}
Without hints, the expected search tries $n/2$ rules on average.
With a perfect hint, the correct rule is tried first.  The
speedup factor is approximately $n/2 / 1 = n/2$.  Under the
convention of worst-case analysis, the exhaustive search tries
all $n$ rules, while the perfect hint tries 1.  Therefore
$E = n$.
\end{proof}

\subsection{Calibration Error Bounds}

\begin{definition}[Expected calibration error]%
\label{def:ece}
The \emph{expected calibration error} (ECE) is:
\[
  \mathrm{ECE} = \sum_{b=1}^{B} \frac{n_b}{N}
  |p_b - c_b|
\]
where $B$ is the number of calibration buckets, $n_b$ is the
count in bucket~$b$, $N$ is the total count, $p_b$ is the
observed accuracy, and $c_b$ is the average confidence.
\end{definition}

\begin{lemma}[ECE bound after $N$ observations]%
\label{lem:ecebound}
After $N$ observations, the ECE is bounded by
$O(B/\sqrt{N_{min}})$ where $N_{min} = \min_b n_b$
is the smallest bucket count.
\end{lemma}
\begin{proof}
In each bucket, the observed accuracy $p_b$ deviates from
the true accuracy by $O(1/\sqrt{n_b})$ with high probability
(by Hoeffding's inequality).  The ECE is a weighted sum of
these deviations.  With $B$ buckets and minimum count
$N_{min}$, the ECE is bounded by $B / \sqrt{N_{min}}$.
\end{proof}

\begin{lemma}[Calibration improves with data]%
\label{lem:calibimprove}
The ECE is non-increasing in expectation as $N$ grows.
\end{lemma}
\begin{proof}
Each new observation refines the empirical accuracy in its
bucket.  By the law of large numbers, the empirical accuracy
converges to the true accuracy.  Isotonic regression then
produces a better fit.  Therefore the ECE decreases.
\end{proof}

\subsection{Session and Rate-Limit Properties}

\begin{lemma}[Deduction sessions are isolated]%
\label{lem:dedsessiso}
Concurrent deduction sessions do not share mutable state.
\end{lemma}
\begin{proof}
Each \texttt{APISession} instance has private fields.  The
\texttt{APIEventLog} is shared but append-only.  Therefore
sessions do not interfere.
\end{proof}

\begin{lemma}[Rate limit compliance for deduction]%
\label{lem:dedratelimit}
The deduction assistant respects the global rate limit.
\end{lemma}
\begin{proof}
All API calls go through \texttt{JuGeoAPI}, which enforces
the rate limit via \texttt{APIRateLimiter}.  The assistant
does not bypass the API.
\end{proof}

\begin{proposition}[Deduction pipeline safety]%
\label{prop:dedpipelinesafe}
The deduction pipeline (hint generation, validation, application)
is safe: only correct rule applications are applied, trust is
bounded by $\Tcopilot$ for hints and $\Tsolver$ for validations,
and the pipeline terminates.
\end{proposition}
\begin{proof}
Correctness: Theorem~\ref{thm:hintsound} and
Proposition~\ref{prop:e2eded}.
Trust: hints are at $\Tcopilot$; validations at $\Tsolver$.
Termination: Lemma~\ref{lem:hintsearchterm}.
\end{proof}

\end{document}
